% GENERATED FILE. Edit canonical lessons, then run npm run generate:books.
% canonical-source-hash: fnv1a64:e5ef26b5539934a3
% canonical-lessons: LA-C506-comis, LA-C506-comitas, LA-C506-urbanitas, LA-C506-rusticitas, LA-C506-liberalitas

\chapter{Manners}
\label{ch:la-manners}
\hlchaptermodality{506}

\begin{chapteropening}
\textbf{By the end of this chapter:} \emph{I can describe good and bad manners.}

Complete the last lesson of chapter 506: I can describe good and bad manners.
\end{chapteropening}

\section[cōmis, cōme]{cōmis, cōme — courteous, friendly}

\label{lesson:LA-C506-comis}



\begin{quote}
Before the new one: say the Latin for kindly, then the Latin for goodwill.
\end{quote}

\subsection*{You'll want to know: cōmis, cōme}
\textbf{cōmis, cōme} — \textquotedblleft{}courteous, friendly\textquotedblright{}.

\begin{quote}
Hospes cōmis omnēs benignē salūtat.
\end{quote}

\begin{quote}
\emph{The courteous host greets everyone kindly.}
\end{quote}

Two endings. Warm, easy, approachable manners.

\begin{cousinweb}
\textbf{cōmis} comes from Old Latin \emph{cosmis} (EtymDB), which EtymDB traces to a Proto-Indo-European root meaning \textquotedblleft{}to smile\textquotedblright{}; Lewis \& Short calls it doubtful, so the origin is not fully settled.
\end{cousinweb}

\begin{grammarlens}[title={What you've built}]
One more word for telling, explaining and describing at B1.
\end{grammarlens}

\subsection*{Your turn}
\begin{itemize}
\raggedright
  \item \emph{Say it:} \emph{cōmis, cōme}
  \item \emph{Say it:} \emph{cōmis, cōme}, once more
  \item \emph{Say it:} a sentence of your own with \emph{cōmis, cōme}
  \item \emph{Recall:} say the Latin for kindly, then the Latin for goodwill, then say \emph{cōmis, cōme} again
\end{itemize}

\begin{tcolorbox}[breakable,colback=teal!4,colframe=teal!35!black,title={Before you move on}]
What does cōmis, cōme mean? (\textquotedblleft{}Courteous, friendly\textquotedblright{}.) Say it once more.
\end{tcolorbox}
\section[cōmitās, cōmitātis]{cōmitās, cōmitātis — friendliness, affability}

\label{lesson:LA-C506-comitas}



\begin{quote}
Before the new one: say the Latin for goodwill, then the Latin for courteous.
\end{quote}

\subsection*{You'll want to know: cōmitās, cōmitātis}
\textbf{cōmitās, cōmitātis} — \textquotedblleft{}friendliness, affability\textquotedblright{}.

\begin{quote}
Cōmitāte suā multōs amīcōs habet.
\end{quote}

\begin{quote}
\emph{Through his friendliness he has many friends.}
\end{quote}

Feminine, third declension, from \emph{cōmis}.

\begin{cousinweb}
\textbf{cōmitās} is built on \emph{cōmis} \textquotedblleft{}courteous\textquotedblright{} with \emph{-itās} (Lewis \& Short; EtymDB).
\end{cousinweb}

\begin{grammarlens}[title={What you've built}]
One more word for telling, explaining and describing at B1.
\end{grammarlens}

\subsection*{Your turn}
\begin{itemize}
\raggedright
  \item \emph{Say it:} \emph{cōmitās, cōmitātis}
  \item \emph{Say it:} \emph{cōmitās, cōmitātis}, once more
  \item \emph{Say it:} a sentence of your own with \emph{cōmitās, cōmitātis}
  \item \emph{Recall:} say the Latin for goodwill, then the Latin for courteous, then say \emph{cōmitās, cōmitātis} again
\end{itemize}

\begin{tcolorbox}[breakable,colback=teal!4,colframe=teal!35!black,title={Before you move on}]
What does cōmitās, cōmitātis mean? (\textquotedblleft{}Friendliness, affability\textquotedblright{}.) Say it once more.
\end{tcolorbox}
\section[urbānitās, urbānitātis]{urbānitās, urbānitātis — refinement, polished wit}

\label{lesson:LA-C506-urbanitas}



\begin{quote}
Before the new one: say the Latin for courteous, then the Latin for friendliness.
\end{quote}

\subsection*{You'll want to know: urbānitās, urbānitātis}
\textbf{urbānitās, urbānitātis} — \textquotedblleft{}refinement, polished wit\textquotedblright{}.

\begin{quote}
Urbānitās eius in sermōne appāret.
\end{quote}

\begin{quote}
\emph{His refinement shows in his conversation.}
\end{quote}

Feminine, third declension, from \emph{urbānus}: the polish of the city, \emph{urbs}. English \emph{urbanity}.

\begin{cousinweb}
\textbf{urbānitās} is built on \emph{urbānus} \textquotedblleft{}of the city\textquotedblright{}, from \emph{urbs} \textquotedblleft{}city\textquotedblright{} (Lewis \& Short; EtymDB). English \emph{urbanity} comes from it.
\end{cousinweb}

\begin{grammarlens}[title={What you've built}]
One more word for telling, explaining and describing at B1.
\end{grammarlens}

\subsection*{Your turn}
\begin{itemize}
\raggedright
  \item \emph{Say it:} \emph{urbānitās, urbānitātis}
  \item \emph{Say it:} \emph{urbānitās, urbānitātis}, once more
  \item \emph{Say it:} a sentence of your own with \emph{urbānitās, urbānitātis}
  \item \emph{Recall:} say the Latin for courteous, then the Latin for friendliness, then say \emph{urbānitās, urbānitātis} again
\end{itemize}

\begin{tcolorbox}[breakable,colback=teal!4,colframe=teal!35!black,title={Before you move on}]
What does urbānitās, urbānitātis mean? (\textquotedblleft{}Refinement, polished wit\textquotedblright{}.) Say it once more.
\end{tcolorbox}
\section[rūsticitās, rūsticitātis]{rūsticitās, rūsticitātis — country manners, boorishness}

\label{lesson:LA-C506-rusticitas}



\begin{quote}
Before the new one: say the Latin for friendliness, then the Latin for refinement.
\end{quote}

\subsection*{You'll want to know: rūsticitās, rūsticitātis}
\textbf{rūsticitās, rūsticitātis} — \textquotedblleft{}country manners, boorishness\textquotedblright{}.

\begin{quote}
Rūsticitātem eius urbānī rīdent.
\end{quote}

\begin{quote}
\emph{City people laugh at his country manners.}
\end{quote}

Feminine, third declension, from \emph{rūsticus}. The opposite of \emph{urbānitās}. English \emph{rusticity}.

\begin{cousinweb}
\textbf{rūsticitās} is built on \emph{rūsticus} \textquotedblleft{}of the country\textquotedblright{}, from \emph{rūs} \textquotedblleft{}countryside\textquotedblright{} (Lewis \& Short). English \emph{rusticity} comes from it.
\end{cousinweb}

\begin{grammarlens}[title={What you've built}]
One more word for telling, explaining and describing at B1.
\end{grammarlens}

\subsection*{Your turn}
\begin{itemize}
\raggedright
  \item \emph{Say it:} \emph{rūsticitās, rūsticitātis}
  \item \emph{Say it:} \emph{rūsticitās, rūsticitātis}, once more
  \item \emph{Say it:} a sentence of your own with \emph{rūsticitās, rūsticitātis}
  \item \emph{Recall:} say the Latin for friendliness, then the Latin for refinement, then say \emph{rūsticitās, rūsticitātis} again
\end{itemize}

\begin{tcolorbox}[breakable,colback=teal!4,colframe=teal!35!black,title={Before you move on}]
What does rūsticitās, rūsticitātis mean? (\textquotedblleft{}Country manners, boorishness\textquotedblright{}.) Say it once more.
\end{tcolorbox}
\section[līberālitās, līberālitātis]{līberālitās, līberālitātis — generosity}

\label{lesson:LA-C506-liberalitas}



\begin{quote}
Before the new one: say the Latin for refinement, then the Latin for country manners.
\end{quote}

\subsection*{You'll want to know: līberālitās, līberālitātis}
\textbf{līberālitās, līberālitātis} — \textquotedblleft{}generosity\textquotedblright{}.

\begin{quote}
Līberālitāte dominī servus lībertātem accipit.
\end{quote}

\begin{quote}
\emph{Through his master's generosity the slave receives his freedom.}
\end{quote}

Feminine, third declension, from \emph{līberālis} — befitting a free person. English \emph{liberality}.

\begin{cousinweb}
\textbf{līberālitās} is built on \emph{līberālis} \textquotedblleft{}befitting a free person\textquotedblright{}, from \emph{līber} \textquotedblleft{}free\textquotedblright{} (Lewis \& Short). English \emph{liberality} comes from it.
\end{cousinweb}

\begin{grammarlens}[title={What you've built}]
One more word for telling, explaining and describing at B1.
\end{grammarlens}

\subsection*{Your turn}
\begin{itemize}
\raggedright
  \item \emph{Say it:} \emph{līberālitās, līberālitātis}
  \item \emph{Say it:} \emph{līberālitās, līberālitātis}, once more
  \item \emph{Say it:} a sentence of your own with \emph{līberālitās, līberālitātis}
  \item \emph{Recall:} say the Latin for refinement, then the Latin for country manners, then say \emph{līberālitās, līberālitātis} again
\end{itemize}

\begin{tcolorbox}[breakable,colback=teal!4,colframe=teal!35!black,title={Before you move on}]
What does līberālitās, līberālitātis mean? (\textquotedblleft{}Generosity\textquotedblright{}.) Say it once more.
\end{tcolorbox}
